% Pala con profilo NACA 2412 approssimato.
% Uso: \DrawBlade{x}{y}{angolo}


% ============================================================
% PARAMETRI MODIFICABILI
% ============================================================
\def\Denoms{150}

% Triangolo di ingresso
\def\OoneX{0}
\def\OoneY{0}

% Estremo di V1
% VoneX diverso da OoneX rende V1 inclinato.
\def\VoneX{618/\Denoms}
\def\VoneY{-385/\Denoms}

% Estremo di W1
\def\WoneX{298/\Denoms}
\def\WoneY{-385/\Denoms}

% Valori visualizzati
\def\VoneValue{728}
\def\WoneValue{386}
\def\UoneValue{320}

% Angoli visualizzati
\pgfmathsetmacro{\AlphaOneValue}{round(10*atan2(abs(\VoneX),abs(\VoneY)))/10}
\def\BetaOneValue{58.1}

% Triangolo di uscita
\def\OtwoX{9}
\def\OtwoY{-5.7}

% Estremo di V2
\def\VtwoX{-124/\Denoms+\OtwoX}
\def\VtwoY{-385/\Denoms+\OtwoY}

% Estremo di W2
\def\WtwoX{-444/\Denoms+\OtwoX}
\def\WtwoY{-385/\Denoms+\OtwoY}

% Valori visualizzati
\def\VtwoValue{404}
\def\WtwoValue{587}
\def\UtwoValue{320}

% Angoli visualizzati

%\pgfmathsetmacro{\AlphaTwoValue}{round(10*atan2(abs(\VtwoX-\OtwoX),abs(\VtwoY-\OtwoY)))/10}

\def\AlphaTwoValue{17.9}

\def\BetaTwoValue{40}

% Pale
\def\BladeAngle{-105}
\def\BladeY{-3.3}

\begin{figure}[H]
\centering
\begin{tikzpicture}[
  scale=0.8,
  font=\fontsize{7}{8}\selectfont,
  absolute/.style={ultra thick,cyan!80!blue,-{Stealth[length=3mm,width=2.2mm]}},
  relative/.style={ultra thick,orange,-{Stealth[length=3mm,width=2.2mm]}},
  peripheral/.style={ultra thick,green!70!black,-{Stealth[length=3mm,width=2.2mm]}},
  reference/.style={dashed,gray!60,thin},
  velocitylabel/.style={text=black,fill=white,inner sep=2pt},
  angletext/.style={text=black,fill=white,inner sep=1pt}
]

% ============================================================
% TRIANGOLO DI INGRESSO
% ============================================================

\coordinate (O1) at (\OoneX,\OoneY);
\coordinate (V1) at (\VoneX,\VoneY);
\coordinate (W1) at (\WoneX,\WoneY);
\coordinate (VerticalOne) at (\OoneX,-3.1);

% Linea verticale di riferimento
\draw[reference] (O1) -- (VerticalOne);

% Vettore V1
\draw[absolute]
  (O1) -- (V1)
  node[velocitylabel,midway,right=10pt]
  {$v_{1.5}=\VoneValue\,\mathrm{m/s}$};

% Vettore W1
\draw[relative]
  (O1) -- (W1)
  node[velocitylabel,midway,below left=1pt]
  {$w_{1.5}=\WoneValue\,\mathrm{m/s}$};

% Vettore U1
\draw[peripheral]
  (W1) -- (V1)
  node[velocitylabel,midway,below=5pt]
  {$U=\UoneValue\,\mathrm{m/s}$};

% Arco beta 1
%\pic[
%  draw,
%  angle radius=1.25cm
%]
%  {angle=VerticalOne--O1--W1};
%
%\node[angletext]
%  at (0.70,-1.15)
%  {$\beta_1=\BetaOneValue^\circ$};

% Arco alpha 1
% Diventa visibile quando V1 viene inclinato.
\pic[
  draw,
  angle radius=0.72cm
]
  {angle=VerticalOne--O1--V1};

\node[angletext]
  at (1.6,0)
  {$\alpha_{1.5} = \pgfmathprintnumber[fixed, precision=1]{\AlphaOneValue}^\circ$}


% ============================================================
% PALE NACA
% ============================================================

\DrawBlade{0}{\BladeY}{\BladeAngle}
\DrawBlade{2}{\BladeY}{\BladeAngle}
\DrawBlade{4}{\BladeY}{\BladeAngle}
\DrawBlade{6}{\BladeY}{\BladeAngle}
\DrawBlade{8}{\BladeY}{\BladeAngle}




% ============================================================
% FRECCIA DI ROTAZIONE
% ============================================================

\node[
  single arrow,
  draw=green!70!black,
  fill=white,
  thick,
  minimum height=1cm,
  minimum width=0.5cm,
  single arrow head extend=0.1cm,
  anchor=east,
  shape border rotate=0
]
  (rotor) at (-0.5,-4.3) {};

\node[
  left=4pt,
  align=center,
  font=\sffamily\scriptsize
]
  at (rotor.west)
  {\textbf{Rotazione}\\\textbf{del rotore}};


% ============================================================
% PIANO DI USCITA
% ============================================================

\draw[gray!65,thin]
  (-1,-5.45) -- (9,-5.45);

\node[below,text=gray!65,font=\scriptsize]
  at (3.2,-5.45)
  {piano di uscita del rotore};


% ============================================================
% TRIANGOLO DI USCITA
% ============================================================

\coordinate (O2) at (\OtwoX,\OtwoY);
\coordinate (V2) at (\VtwoX,\VtwoY);
\coordinate (W2) at (\WtwoX,\WtwoY);
\coordinate (VerticalTwo) at (\OtwoX,-8.3);

% Linea verticale di riferimento
\draw[reference] (O2) -- (VerticalTwo);

% Vettore V2
\draw[absolute]
  (O2) -- (V2)
  node[velocitylabel, midway, below right=8pt, xshift=-4pt]
  {$v_2=\VtwoValue\,\mathrm{m/s}$};

% Vettore W2
\draw[relative]
  (O2) -- (W2)
  node[velocitylabel,midway,left=10pt]
  {$w_2=\WtwoValue\,\mathrm{m/s}$};

% Vettore U2
\draw[peripheral]
  (W2) -- (V2)
  node[velocitylabel,midway,below=8pt]
  {$U=\UtwoValue\,\mathrm{m/s}$};

% Arco alpha 2
\pic[
  draw,
  angle radius=0.5cm
]
  {angle=V2--O2--VerticalTwo};

\node[angletext]
  at (10.1,-6.3)
  {$\alpha_2=\AlphaTwoValue^\circ$};

% Arco beta 2
%\pic[
%  draw,
%  angle radius=1.55cm
%]
%  {angle=VerticalTwo--O2--W2};
%
%\node[angletext]
%  at (9.80,-6.65)
%  {$\beta_2=\BetaTwoValue^\circ$};
\draw[->, thick] (9,0) -- (10.5, 0) node[right] {$\theta$};
% Asse theta (verticale)
\draw[->, thick] (9,0) -- (9, -1.5) node[below] {$z$};
% Origine
\fill (0,0) circle (1.5pt);

  
\end{tikzpicture}
    \caption{Diagramma dei triangoli di velocità all'ingresso (1.5) e all'uscita (2) del rotore.}
    \label{fig:turb:triangoli}
\end{figure}


